\section{Conclusion}

\subsection{Proof-of-Concept Assessment}

This technical validation establishes a **solid proof-of-concept** for a patient-specific, SOM-based EEG emotion recognition pipeline. We have successfully developed and validated the fundamental technological building blocks:

\begin{itemize}
    \item A practical 8-channel dry-electrode EEG system (setup < 3 minutes)
    \item A real-time acquisition script with robust timestamp synchronization
    \item An intuitive annotation interface for non-expert users
    \item A comprehensive feature extraction module (spectral, asymmetry, connectivity)
    \item An unsupervised SOM framework for patient-specific representation learning
    \item A supervised classification pipeline with interpretable feature importance
\end{itemize}

Validation on a healthy volunteer (subject \subject{F42}, age 42) demonstrates that all components function as intended, with promising preliminary results: 77.8\% classification accuracy across six emotional states with only 4.7 minutes of recording, and identification of connectivity variability and theta power as key discriminative features.

\subsection{From Proof-of-Concept to Application}

This work represents the **first step** toward a non-invasive, easy-to-deploy emotion recognition system. The current validation, while encouraging, was conducted on a healthy volunteer in a controlled environment. The maturation phase would aim to:

\begin{itemize}
    \item Test the system in **ecological conditions** (home, private practice settings)
    \item Evaluate usability by **non-expert users** (families, caregivers, ergotherapists)
    \item Collect data from **diverse populations** to validate the patient-specific approach
    \item Refine the interface based on real-world user feedback
\end{itemize}

\subsection{Potential Applications}

The system's simplicity and non-invasive nature open numerous application avenues:

\begin{itemize}
    \item \textbf{Home use}: Enabling families to better understand emotions of a non-verbal relative
    \item \textbf{Private practice}: Assistive tool for ergotherapists, psychologists, educators
    \item \textbf{Care facilities}: Additional information for support staff
    \item \textbf{Research}: Robust platform for studying emotions in various populations
\end{itemize}

\subsection{Maturation Phase Requirements}

To move from this proof-of-concept to real-world applications, several developments are needed:

\begin{enumerate}
    \item \textbf{System simplification}: Further setup time reduction, interface improvement
    \item \textbf{Long-term robustness}: Extended testing in uncontrolled environments
    \item \textbf{Multi-user validation}: Evaluation with different operators
    \item \textbf{Application-specific tuning}: Feature and model adaptation by use case
    \item \textbf{Transfer preparation}: Documentation, packaging, patent filings
\end{enumerate}

\subsection{Conclusion and Funding Call}

This work establishes a **solid technological foundation** for a patient-specific EEG emotion recognition system. The building blocks are validated, the approach is sound, and potential applications are numerous. What is now needed is support to move from laboratory to field.

CNRS pre-maturation funding would enable us to:
\begin{itemize}
    \item Conduct necessary validation studies in real-world conditions
    \item Refine the system based on user feedback
    \item Prepare technology transfer and broad dissemination
    \item Establish partnerships with potential users and industry partners
    \item Protect intellectual property and valorize innovations
\end{itemize}

The science is ready; the technology works; the next step is to bring it to those who could benefit from it.